\subsection{Implementasi \textit{From Scratch}}
\label{subsec:nb-implementasi}

\subsubsection{Landasan Teori}
\todo{Jelaskan teorema Bayes, asumsi independensi bersyarat (\textit{Naive Bayes} sebagai \textit{Bayesian network} dengan kelas sebagai \textit{parent} semua fitur), \textit{likelihood} kategorikal dengan Laplace \textit{smoothing} dan numerik dengan Gaussian, serta perhitungan dalam ruang log. Acu \cite{slide-bayesnet}. Sertakan rumus utama dalam lingkungan \texttt{equation}/\texttt{align}.}

\subsubsection{Struktur Kelas dan Fungsi}
\todo{Deskripsikan kelas pada \texttt{src/algoritma/nb/nb.py}: \texttt{\_\_init\_\_}, \texttt{fit} (estimasi \textit{prior} dan \textit{likelihood}), \texttt{predict\_proba}, \texttt{predict}, \texttt{save}, \texttt{load}. Isi tabel di bawah, satu baris per fungsi/metode.}

\begin{table}[H]
	\centering
	\caption{Deskripsi kelas dan fungsi NB}
	\label{tab:nb-fungsi}
	\begin{tabularx}{\textwidth}{l Y}
		\toprule
		\textbf{Fungsi/Metode} & \textbf{Deskripsi} \\
		\midrule
		\texttt{nama\_fungsi} & \todo{deskripsi} \\
		\bottomrule
	\end{tabularx}
\end{table}

\subsubsection{\textit{Source Code}}
\todo{Sisipkan potongan kode inti (estimasi parameter dan perhitungan \textit{posterior}) dengan makro \texttt{\textbackslash kode}, lalu jelaskan alurnya. Cek ulang nomor baris jika kode berubah.}
% \kode{../../src/algoritma/nb/nb.py}{1}{30}{Judul potongan kode}{lst:nb-inti}

\subsubsection{\textit{Save} dan \textit{Load} Model}
\todo{Format penyimpanan (JSON/\texttt{.npz}/\textit{pickle}), isi yang disimpan, dan contoh pemakaian.}
